\section*{Singular values and orthogonal coordinates}
All matrices in this unit are real. A matrix \(A\) of size \(m\times n\)
may be rectangular or singular. The symmetric spectral theorem applies
to its Gram matrices \(A^\top A\) and \(AA^\top\), since
\(x^\top A^\top Ax=\norm{Ax}^2\geqslant0\). Moreover,
\[
\ker(A^\top A)=\ker A:
\quad A^\top Ax=0\ \Longrightarrow\ \norm{Ax}^2=0\ \Longrightarrow\ Ax=0,
\]
and the reverse inclusion follows by multiplication. Thus
\(\operatorname{rank}(A^\top A)=\operatorname{rank}A\). Apply the same
argument to \(A^\top\) for \(AA^\top\).

\begin{theorem}[Singular value decomposition]
Let \(A\) have rank \(r>0\). There are orthonormal lists
\(u_1,\ldots,u_r\in\R^m\), \(v_1,\ldots,v_r\in\R^n\), and numbers
\(\sigma_1\geqslant\cdots\geqslant\sigma_r>0\) such that
\[
Av_i=\sigma_i u_i,\qquad A^\top u_i=\sigma_i v_i,
\qquad A=\sum_{i=1}^r\sigma_i u_i v_i^\top.
\]
Writing \(U_r=(u_1,\ldots,u_r)\), \(V_r=(v_1,\ldots,v_r)\), and
\(D_r=\operatorname{diag}(\sigma_1,\ldots,\sigma_r)\), this is
\(A=U_rD_rV_r^\top\), with \(U_r^\top U_r=V_r^\top V_r=I_r\).
Extending both lists to orthonormal bases gives \(A=U\Sigma V^\top\),
where \(U,V\) are square orthogonal matrices and the only nonzero entries
of the \(m\times n\) matrix \(\Sigma\) are its first \(r\) diagonal entries
\(\sigma_i\). For \(A=0\), use \(\Sigma=0\) and any orthogonal \(U,V\).
\end{theorem}
\begin{proof}
Choose an orthonormal eigenbasis \(v_1,\ldots,v_n\) of \(A^\top A\),
putting its \(r\) positive eigenvalues \(\lambda_i\) first, in decreasing
order. Set \(\sigma_i=\sqrt{\lambda_i}\) and \(u_i=Av_i/\sigma_i\).
Then
\[
u_i^\top u_j=\frac{v_i^\top A^\top Av_j}{\sigma_i\sigma_j}
=\frac{\lambda_jv_i^\top v_j}{\sigma_i\sigma_j}
=\begin{cases}1&i=j,\\0&i\ne j.\end{cases}
\]
Also \(A^\top u_i=\lambda_i v_i/\sigma_i=\sigma_i v_i\).
For \(i>r\), \(v_i\in\ker(A^\top A)=\ker A\), so \(Av_i=0\).
Expanding an arbitrary \(x\) in the \(v_i\) basis gives
\[
Ax=\sum_{i=1}^r\sigma_i u_i(v_i^\top x),
\]
proving the matrix identity. Extend the \(u_i\) to an orthonormal basis
of \(\R^m\). The identities for \(Av_i\) say exactly that
\(U^\top AV=\Sigma\), hence \(A=U\Sigma V^\top\).
\end{proof}

The \(\sigma_i\) are the positive \textbf{singular values} of \(A\).
When all singular values are listed, append zeros to make
\(\min(m,n)\) entries. The eigenvalues of \(A^\top A\) include these
squares and, when necessary, further zeros to make \(n\) entries.
The positive eigenvalues of \(A^\top A\) and \(AA^\top\) agree, including
multiplicities: the constructed \(u_i\) are orthonormal eigenvectors of
\(AA^\top\), and its rank is also \(r\).

The factors on the two sides are linked. Choose eigenvectors of
\(A^\top A\), then define \(u_i=Av_i/\sigma_i\). Independently choosing
signs or bases in both Gram eigenspaces can break that identity.
For example, \(A=(1)\), \(u_1=1\), \(v_1=-1\) satisfy both Gram
eigenvalue equations but give \(u_1\sigma_1v_1^\top=-1\), not \(A\).
When a positive singular value repeats, the same orthogonal change must
be made in its left and right singular-vector lists.

\begin{example}[A rectangular decomposition]\label{eg:10:rectangular}
Let
\[
A=\begin{pmatrix}1&0&1\\0&1&0\end{pmatrix},\qquad
A^\top A=\begin{pmatrix}1&0&1\\0&1&0\\1&0&1\end{pmatrix}.
\]
In the plane of first and third coordinates, \((1,0,1)^\top\) and
\((-1,0,1)^\top\) have eigenvalues \(2,0\); \(e_2\) has eigenvalue \(1\).
Normalize in the order \(2,1,0\):
\[
v_1=\frac1{\sqrt2}(1,0,1)^\top,\qquad
v_2=(0,1,0)^\top,\qquad
v_3=\frac1{\sqrt2}(-1,0,1)^\top.
\]
The positive singular values are \(\sqrt2,1\), and
\[
u_1=Av_1/\sqrt2=(1,0)^\top,\qquad
u_2=Av_2=(0,1)^\top.
\]
Thus \(U=I_2\),
\[
V=\begin{pmatrix}1/\sqrt2&0&-1/\sqrt2\\0&1&0\\1/\sqrt2&0&1/\sqrt2\end{pmatrix},
\qquad \Sigma=\begin{pmatrix}\sqrt2&0&0\\0&1&0\end{pmatrix},
\qquad U\Sigma V^\top=A.
\]
The first row of \(\Sigma V^\top\) is \((1,0,1)\), and the second is
\((0,1,0)\), verifying reconstruction. The formula
\(Ax=\sqrt2u_1(v_1^\top x)+u_2(v_2^\top x)\)
shows which direction is lost and which directions are stretched.
\end{example}

\subsection*{The four subspaces and the size of a matrix}
The SVD gives orthonormal bases for the four fundamental subspaces:
\[
\begin{array}{ll}
\operatorname{im}A=\operatorname{span}(u_1,\ldots,u_r),&
\ker A^\top=\operatorname{span}(u_{r+1},\ldots,u_m),\\
\operatorname{im}A^\top=\operatorname{span}(v_1,\ldots,v_r),&
\ker A=\operatorname{span}(v_{r+1},\ldots,v_n).
\end{array}
\]
The expansion of \(Ax\) gives the image inclusion, and
\(u_i=A(v_i/\sigma_i)\) gives its reverse. Independence of the \(u_i\)
makes \(Ax=0\) equivalent to \(v_i^\top x=0\) for \(i\leqslant r\).
Apply the same argument to \(A^\top\). In particular \(A\) restricts to
a bijection \(\operatorname{im}A^\top\to\operatorname{im}A\).

For \(x=\sum_i z_i v_i\), orthogonality gives
\(\norm{Ax}^2=\sum_{i=1}^r\sigma_i^2z_i^2\).
Define \(\norm{A}_2=\max_{\norm{x}=1}\norm{Ax}\), the operator norm.
For \(r>0\), this identity gives \(\norm{A}_2=\sigma_1\), attained at
\(v_1\). On \(\operatorname{im}A^\top\), it also gives
\[
\sigma_r\norm{x}\leqslant\norm{Ax}\leqslant\sigma_1\norm{x}.
\]
The Frobenius norm satisfies
\[
\norm{A}_F^2=\tr(A^\top A)=\sum_{i=1}^r\sigma_i^2.
\]
Indeed \(A^\top A=V_rD_r^2V_r^\top\) and cyclicity of trace give
\(\tr(A^\top A)=\tr(D_r^2V_r^\top V_r)=\tr(D_r^2)\).

\section*{The Moore--Penrose inverse}
The preceding bijection has an inverse on the image. Extend it by zero
on the perpendicular complement of the image. Its matrix is
\[
A^+=V_rD_r^{-1}U_r^\top
=\sum_{i=1}^r\frac1{\sigma_i}v_i u_i^\top,
\qquad 0^+=0_{n\times m}.
\]
This \textbf{Moore--Penrose inverse} has size \(n\times m\).
For nonsingular square \(A\), it is the ordinary inverse.
For all \(A\), multiplication gives
\[
AA^+=U_rU_r^\top=:P,\qquad A^+A=V_rV_r^\top=:Q.
\]
Thus \(P\) projects orthogonally onto \(\operatorname{im}A\), and \(Q\)
onto \(\operatorname{im}A^\top\). They generally differ, possibly in size.

\begin{theorem}[Four equations and uniqueness]
The matrix \(A^+\) is the unique \(n\times m\) matrix \(X\) satisfying
\[
AXA=A,\qquad XAX=X,\qquad (AX)^\top=AX,\qquad (XA)^\top=XA.
\]
Consequently its definition is independent of the chosen SVD.
\end{theorem}
\begin{proof}
The SVD formula gives symmetric \(AA^+=P\), \(A^+A=Q\), and
\[
AA^+A=U_rU_r^\top U_rD_rV_r^\top=A,\qquad
A^+AA^+=V_rV_r^\top V_rD_r^{-1}U_r^\top=A^+.
\]
For uniqueness, suppose \(X\) satisfies the four equations.
The matrix \(AX\) is symmetric idempotent, and its image is
\(\operatorname{im}A\): one inclusion follows from its factor \(A\),
and \(AX(Az)=Az\) gives the reverse. Hence \(AX=P\).
Likewise \(XA\) is symmetric idempotent and
\(\ker(XA)=\ker A\), since \(XAz=0\) implies \(Az=AXAz=0\).
Its image is therefore \((\ker A)^\perp=\operatorname{im}A^\top\), so
\(XA=Q\). Finally
\[
X=XAX=QX=A^+AX=A^+P=A^+.
\]
Here uniqueness of the orthogonal projector follows from uniqueness of
the orthogonal decomposition.
\end{proof}

\begin{example}[Recovering the part that can be recovered]
For the rectangular matrix of Example~\ref{eg:10:rectangular},
\[
A^+=\frac1{\sqrt2}v_1u_1^\top+v_2u_2^\top
=\begin{pmatrix}1/2&0\\0&1\\1/2&0\end{pmatrix},
\quad AA^+=I_2,\quad
A^+A=\begin{pmatrix}1/2&0&1/2\\0&1&0\\1/2&0&1/2\end{pmatrix}.
\]
For \(b=(2,3)^\top\), \(x_0=A^+b=(1,3,1)^\top\) satisfies \(Ax_0=b\).
All exact solutions are
\[
x=(1,3,1)^\top+t(-1,0,1)^\top,\qquad t\in\R,
\]
since \((-1,0,1)^\top\) spans the kernel. The two terms are orthogonal,
so \(\norm{x}^2=11+2t^2\). The inverse selects the unique solution of
smallest norm, though there are infinitely many exact solutions.
\end{example}

\section*{Least squares as orthogonal projection}
Given \(b\in\R^m\), least squares asks for a vector \(x\) minimizing
\(\norm{Ax-b}^2\). Set
\[
p=Pb=AA^+b,\qquad r=b-p,\qquad x_0=A^+b.
\]
Here \(p\in\operatorname{im}A\), \(r\perp\operatorname{im}A\), and
\(Ax_0=p\). For every \(x\), Pythagoras therefore gives
\[
Ax-b=A(x-x_0)-r,\qquad
\norm{Ax-b}^2=\norm{A(x-x_0)}^2+\norm{r}^2.
\]
This proves existence without assuming full rank or consistency.

\begin{theorem}[All least-squares solutions]
For any real \(m\times n\) matrix \(A\) and any \(b\in\R^m\),
the following are equivalent:
\begin{enumerate}
\item \(x\) minimizes \(\norm{Ax-b}^2\);
\item \(Ax=AA^+b\);
\item \(A^\top Ax=A^\top b\), the \textbf{normal equations};
\item \(x=A^+b+z\) for some \(z\in\ker A\), equivalently
\(x=A^+b+(I_n-A^+A)w\) for arbitrary \(w\in\R^n\).
\end{enumerate}
All minimizers give the same fitted vector \(p\) and residual \(r\).
The unique minimizer of smallest norm is \(x_0=A^+b\).
The exact system \(Ax=b\) is consistent if and only if \(Pb=b\).
\end{theorem}
\begin{proof}
The squared-distance identity proves (a)\(\Leftrightarrow\)(b) and
(b)\(\Leftrightarrow\)\(A(x-x_0)=0\), which gives the first form in (d).
The projector \(I-Q\) has image \(\ker A\), giving the second form.
Condition (b) says \(b-Ax\perp\operatorname{im}A\), equivalent to
\(A^\top(b-Ax)=0\), which is (c). Conversely (c) gives that same
orthogonality, hence the unique projection, proving (b).
Since \(x_0\in\operatorname{im}A^\top=(\ker A)^\perp\),
\[
\norm{x_0+z}^2=\norm{x_0}^2+\norm{z}^2\quad(z\in\ker A),
\]
with equality to \(\norm{x_0}^2\) only for \(z=0\).
Finally \(Ax=b\) is solvable exactly when \(b\in\operatorname{im}A\),
exactly when its projection is itself.
\end{proof}

If \(A\) has full column rank, \(A^\top A\) is positive definite because
\(\norm{Ax}^2>0\) for nonzero \(x\). The normal equations then give
\[
x=(A^\top A)^{-1}A^\top b,\qquad A^+=(A^\top A)^{-1}A^\top.
\]
Applying this to \(A^\top\) when \(A\) has full row rank gives
\(A^+=A^\top(AA^\top)^{-1}\). Then \(AA^+=I_m\), so every right-hand
side has an exact solution. Without full column rank, uniqueness of the
fitted vector differs from uniqueness of its coefficients.

\begin{example}[An inconsistent system with dependent columns]\label{eg:10:dependent}
Let
\[
A=\begin{pmatrix}1&2\\1&2\end{pmatrix},\qquad b=\begin{pmatrix}1\\2\end{pmatrix}.
\]
Both rows of \(Ax\) are equal, so the system is inconsistent. Compute
\[
A^\top A=\begin{pmatrix}2&4\\4&8\end{pmatrix},
\quad v_1=\frac1{\sqrt5}(1,2)^\top,
\quad v_2=\frac1{\sqrt5}(-2,1)^\top.
\]
Their eigenvalues are \(10,0\). Thus \(\sigma_1=\sqrt{10}\),
\(u_1=Av_1/\sqrt{10}=(1,1)^\top/\sqrt2\); take
\(u_2=(-1,1)^\top/\sqrt2\). It follows that
\[
A=\sqrt{10}\,u_1v_1^\top,
\quad A^+=\frac1{10}\begin{pmatrix}1&1\\2&2\end{pmatrix},
\quad P=\frac12\begin{pmatrix}1&1\\1&1\end{pmatrix},
\quad Q=\frac15\begin{pmatrix}1&2\\2&4\end{pmatrix}.
\]
Multiplying gives
\[
x_0=\begin{pmatrix}3/10\\3/5\end{pmatrix},\qquad
p=\begin{pmatrix}3/2\\3/2\end{pmatrix},\qquad
r=\begin{pmatrix}-1/2\\1/2\end{pmatrix}.
\]
The normal equations are \(2x_1+4x_2=3\), \(4x_1+8x_2=6\), so all
least-squares solutions are
\[
x=\begin{pmatrix}3/10\\3/5\end{pmatrix}
+t\begin{pmatrix}-2\\1\end{pmatrix},\qquad t\in\R.
\]
Each fits \(p\) and has residual sum of squares \(1/2\).
Its squared norm is \(9/20+5t^2\), showing the minimum-norm choice.
For arbitrary coefficients the entire objective is
\[
\norm{Ax-b}^2=(x_1+2x_2-1)^2+(x_1+2x_2-2)^2
=2(x_1+2x_2-3/2)^2+1/2.
\]
Thus projection, SVD, and the normal equations agree.
\end{example}

\subsection*{Sensitivity and small singular values}
For fixed nonzero \(A\), changing the right-hand side by \(h\) changes
its minimum-norm least-squares solution by
\[
A^+h=\sum_{i=1}^r\frac{u_i^\top h}{\sigma_i}v_i.
\]
Bessel's inequality gives
\[
\norm{A^+h}^2=\sum_{i=1}^r\frac{(u_i^\top h)^2}{\sigma_i^2}
\leqslant\frac1{\sigma_r^2}\norm{h}^2.
\]
Equality holds for \(h=u_r\), so \(\norm{A^+}_2=1/\sigma_r\).
A small positive singular value can strongly amplify one component of
an error in \(b\). A component in \(\ker A^\top\) changes the residual
and has no effect on \(A^+b\).
For full column rank, the eigenvalues of \(A^\top A\) range from
\(\sigma_n^2\) to \(\sigma_1^2\); their ratio is the square of
\(\sigma_1/\sigma_n\). This explains why small singular values matter
when calculating a solution, even though the normal equations give an
exact theoretical characterization.
